\documentclass{article}
\usepackage{amsmath}
\usepackage{xcolor}
\usepackage{mystyletau} 
\usepackage[margin=2.5cm]{geometry} 
\usepackage{hyperref} 
\hypersetup{
    colorlinks=true,
    linkcolor=blue,
    urlcolor=blue,
    pdftitle={The mystyletau package},
}

\usepackage{titling}
\setlength{\droptitle}{-1.5cm} 

\title{The \texttt{mystyletau} package}
\author{\href{mailto:antonellomeccariello234@gmail.com}{Antonello Meccariello} 
  \& 
  \href{mailto:meccariellofrancesco87@gmail.com}{Francesco Meccariello}%
}
\date{Version 2.0 --- April 2026}

\begin{document}
\maketitle

\section{Introduction}
The \texttt{mystyletau} package introduces two custom glyphs for the letter $\tau$. While the standard notation for a topological space $(X, \tau_X)$ is functional, the authors believe these variants offer superior aesthetic appeal for denoting topologies or algebraic structures.

\section{New in Version 2.0}
Version 2.0 marks a significant technical shift. The symbols are no longer external PDF images; they are now implemented as true vector glyphs within a dedicated TrueType font (\texttt{mytautwo.ttf}). This ensures:
\begin{itemize}
    \item \textbf{Mathematical Integration:} Symbols scale perfectly in subscripts ($\vtau_{\vtau}$) and superscripts.
    \item \textbf{Color Support:} Symbols dynamically respond to color commands, e.g., {\color{red}\vtau} or {\color{blue}\atau}.
    \item \textbf{Mode Independence:} Commands work in both text and math mode via \texttt{\textbackslash ensuremath}.
\end{itemize}

\section{Usage}

\begin{quote}
\textbf{\color{red}{Important:}} This package requires the use of \textbf{LuaLaTeX} or \textbf{XeLaTeX}. Compilation with \texttt{pdflatex} will result in a fatal error, as the \texttt{fontspec} package is used to handle the custom vector font.
\end{quote}

The following commands provide access to the custom glyphs:
\begin{itemize}
    \item \verb|\vtau|: \vtau
    \item \verb|\atau|: \atau
\end{itemize}

\section{Examples}
Consider the following statements to compare the visual impact:
\begin{enumerate}
    \item Standard: Let $(X, \tau_X)$ be a topological space.
    \item Variant A: Let $(X, \vtau_X)$ be a topological space.
    \item Variant B: Let $(X, \atau_X)$ be a topological space.
\end{enumerate}

\noindent The authors suggest that notations like $(X, \vtau_X)$ and $(X, \atau_X)$ provide a more elegant visual impact compared to the standard Computer Modern $(X, \tau_X)$.

\end{document}